% Part VI -- Algorithms, Optimization, and Machine Learning
% Chapter 18: Machine learning foundations

\h{Machine Learning Foundations}
\label{ch:18}

Machine learning uses data to estimate a rule that can make predictions or
organize new observations. The foundational challenge is not calling
\c{fit()}. It is defining a legitimate target, protecting evaluation data,
choosing a baseline, measuring error, checking residuals, and limiting claims
to the population and conditions represented by the data.

\begin{NoteBox}
\textbf{Prerequisites.} You should understand functions, NumPy arrays, pandas
DataFrames, plots, and algorithm contracts. This chapter uses reproducible
synthetic building data so the computational workflow can be studied without
privacy or licensing concerns.
\end{NoteBox}

\begin{tcolorbox}[title=Learning outcomes,colframe=orange,colback=white,arc=2mm]
After completing this chapter, you should be able to:
\begin{itemize}
    \item distinguish regression, classification, and clustering;
    \item identify samples, features $X$, and target $y$;
    \item split training and test data reproducibly;
    \item explain why preprocessing and feature decisions must not use test answers;
    \item compare a fitted model with a simple baseline;
    \item fit and predict with \c{LinearRegression};
    \item interpret mean absolute error, $R^2$, and residual plots;
    \item recognize overfitting, leakage, and distribution shift; and
    \item document the limits and intended use of a predictive model.
\end{itemize}
\end{tcolorbox}

\hh{Choose the learning task from the output}

The type of answer determines the task.

\begin{center}
\begin{tabular}{lll}
\textbf{Task} & \textbf{Output} & \textbf{Example} \\
Regression & numerical value & predict annual energy use \\
Classification & category & flag likely comfort complaint \\
Clustering & group without target labels & explore similar operating profiles \\
\end{tabular}
\end{center}

Regression and classification are \emph{supervised}: training rows include a
target answer. Clustering is \emph{unsupervised}: an algorithm groups rows from
feature similarity, but the resulting groups do not automatically have a
meaningful real-world interpretation.

\begin{ExplainBox}{A learned model is an estimated mapping}
Training selects parameter values for a mapping from features $X$ to an output.
The mapping summarizes patterns in the training sample under a chosen model
family and loss function. It does not store scientific truth, establish cause,
or guarantee performance on a different population. Independent test data
estimate how the completed workflow generalizes under conditions represented
by that test sample.
\end{ExplainBox}

\begin{SyntaxBox}{Samples are rows; features and targets have different roles}
\begin{itemize}
    \item One \textbf{sample} is one observation, usually one DataFrame row.
    \item A \textbf{feature} is an input variable available when a prediction is made.
    \item The \textbf{target} is the outcome the supervised model learns to predict.
    \item $X$ conventionally stores a two-dimensional feature table.
    \item $y$ conventionally stores a one-dimensional target sequence.
\end{itemize}
A variable measured after the target outcome occurs is usually not a legitimate
feature for an earlier prediction.
\end{SyntaxBox}

\hh{Install and import scikit-learn}

Install scikit-learn in the environment that runs the script.

\begin{verbatim}
python -m pip install scikit-learn
\end{verbatim}

Import estimators and evaluation tools from their focused modules.

\begin{lstlisting}[
    language=python,
    style=block,
    caption={Import the regression workflow components},
    label={c18_imports},
    captionpos=b
]
from sklearn.linear_model import LinearRegression
from sklearn.metrics import mean_absolute_error, r2_score
from sklearn.model_selection import train_test_split
\end{lstlisting}

The official documentation describes
\href{https://scikit-learn.org/stable/modules/generated/sklearn.model_selection.train_test_split.html}{train-test splitting},
\href{https://scikit-learn.org/stable/modules/generated/sklearn.linear_model.LinearRegression.html}{linear regression},
and
\href{https://scikit-learn.org/stable/modules/generated/sklearn.metrics.mean_absolute_error.html}{mean absolute error}.

\hh{Protect the final test before fitting}

The test set simulates unseen data. If its target values influence cleaning,
feature selection, model choice, or tuning, the final evaluation is no longer
independent.

\bookfigure[0.98\textwidth]{Figures/Generated/ch18_ml_workflow.pdf}
{A defensible workflow defines features and target, protects a test set, fits only on training data, predicts unseen test rows, and evaluates against a baseline.}
{fig:ch18-workflow}

\begin{lstlisting}[
    language=python,
    style=block,
    caption={Split features and target reproducibly},
    label={c18_split},
    captionpos=b
]
X_train, X_test, y_train, y_test = train_test_split(
    X,
    y,
    test_size=0.25,
    random_state=42,
)
\end{lstlisting}

\begin{SyntaxBox}{One call returns four aligned objects}
\c{test_size=0.25} assigns one quarter of the rows to testing. A fixed
\c{random_state} makes the random split reproducible for study and debugging.
Feature rows and target values are shuffled together, preserving alignment.
Fit transformations and models on \c{X_train} and \c{y_train}; use the test
objects only for final prediction and evaluation.
\end{SyntaxBox}

For time-ordered or grouped observations, random row splitting may leak future
or related information. Choose a splitting algorithm that reflects how the
model will encounter new data.

\hh{Establish a baseline before fitting a model}

A model should improve on a simple, relevant rule. For regression, predicting
the training-target mean for every test row is a useful starting baseline.

\begin{lstlisting}[
    language=python,
    style=block,
    caption={Evaluate a training-mean baseline},
    label={c18_baseline},
    captionpos=b
]
import numpy as np

baseline_predictions = np.full(
    len(y_test),
    y_train.mean(),
)
baseline_mae = mean_absolute_error(
    y_test,
    baseline_predictions,
)
print("Baseline MAE:", baseline_mae)
\end{lstlisting}

The baseline uses only \c{y_train} to determine its constant prediction. Using
the test-target mean would leak test answers into the benchmark.

\hh{Fit parameters on training data}

Linear regression estimates an intercept and one coefficient per feature. Its
prediction is a weighted sum:

\[
\hat{y} = b_0 + b_1x_1 + b_2x_2 + \cdots + b_px_p.
\]

The fitting algorithm chooses coefficients that minimize the sum of squared
training residuals.

\begin{lstlisting}[
    language=python,
    style=block,
    caption={Fit a linear model and predict unseen rows},
    label={c18_fit_predict},
    captionpos=b
]
model = LinearRegression()
model.fit(X_train, y_train)

predictions = model.predict(X_test)

print("Intercept:", model.intercept_)
print("Coefficients:", model.coef_)
\end{lstlisting}

\begin{SyntaxBox}{Estimator objects separate configuration, fitting, and use}
\begin{enumerate}
    \item \c{LinearRegression()} constructs an unfitted estimator.
    \item \c{fit(X_train, y_train)} learns parameters and stores them on the object.
    \item Names ending in an underscore, such as \c{coef_}, are learned attributes.
    \item \c{predict(X_test)} applies the learned rule without changing the test targets.
\end{enumerate}
The test feature table must contain the same variables with the same meaning
and units as the training feature table.
\end{SyntaxBox}

\hh{Evaluate magnitude and explanatory fit}

No single metric describes every failure. This chapter pairs mean absolute
error with $R^2$ and a residual plot.

\begin{lstlisting}[
    language=python,
    style=block,
    caption={Calculate two complementary regression metrics},
    label={c18_metrics},
    captionpos=b
]
mae = mean_absolute_error(y_test, predictions)
r2 = r2_score(y_test, predictions)

print(f"Model MAE: {mae:.2f}")
print(f"Test R2: {r2:.3f}")
\end{lstlisting}

Mean absolute error is the average absolute prediction error in the target's
units. Lower is better; zero is perfect. $R^2$ compares squared error with a
constant-mean reference: 1 is perfect, 0 is comparable to that reference, and a
negative value is worse on the evaluated data. Neither metric proves causality
or fairness.

\hh{Inspect residuals rather than trusting one score}

A residual is actual minus predicted. A useful residual plot has points
scattered around zero without a strong curve, funnel, or separated group. A
pattern can indicate nonlinearity, changing variance, omitted features, data
errors, or distinct populations.

\bookfigure[0.96\textwidth]{Figures/Generated/ch18_regression_diagnostics.pdf}
{On protected test data, predicted values follow the equality line and residuals are centered near zero; the displayed metrics summarize but do not replace visual diagnosis.}
{fig:ch18-diagnostics}

The synthetic example performs well because its target was deliberately
generated from a mostly linear relationship. Real building data contain sensor
errors, schedules, weather dependence, interactions, and operational changes.

\hh{Applied example: building-energy regression}
\label{sec:c18-application}

The runnable program \c{examples/c18_building_energy_ml.py} generates documented
synthetic data, protects a test set, calculates a baseline, fits linear
regression, prints coefficients, evaluates metrics, plots residuals, and exports
the diagnostic figure.

\begin{lstlisting}[
    language=python,
    style=block,
    caption={Build and evaluate a reproducible regression workflow},
    label={c18_application},
    captionpos=b
]
import numpy as np
import pandas as pd
from sklearn.linear_model import LinearRegression
from sklearn.metrics import mean_absolute_error, r2_score
from sklearn.model_selection import train_test_split


def make_building_data(seed=42, sample_count=240):
    rng = np.random.default_rng(seed)
    X = pd.DataFrame({
        "floor_area_m2": rng.uniform(80, 520, sample_count),
        "glazing_ratio": rng.uniform(0.12, 0.65, sample_count),
        "occupants": rng.integers(2, 45, sample_count),
        "outdoor_temp_c": rng.uniform(-5, 31, sample_count),
    })
    noise = rng.normal(0, 10, sample_count)
    y = (
        28
        + 0.24 * X["floor_area_m2"]
        + 66 * X["glazing_ratio"]
        + 1.35 * X["occupants"]
        - 2.1 * X["outdoor_temp_c"]
        + noise
    )
    return X, y


X, y = make_building_data()
X_train, X_test, y_train, y_test = train_test_split(
    X,
    y,
    test_size=0.25,
    random_state=42,
)

baseline = np.full(len(y_test), y_train.mean())
baseline_mae = mean_absolute_error(y_test, baseline)

model = LinearRegression()
model.fit(X_train, y_train)
predictions = model.predict(X_test)

mae = mean_absolute_error(y_test, predictions)
r2 = r2_score(y_test, predictions)

print(f"Baseline MAE: {baseline_mae:.2f}")
print(f"Model MAE: {mae:.2f}")
print(f"Test R2: {r2:.3f}")
print(pd.Series(model.coef_, index=X.columns))
\end{lstlisting}

\begin{SyntaxBox}{Trace the application by responsibility}
\begin{enumerate}
    \item The generator makes provenance, random seed, features, target rule,
          and noise explicit.
    \item The split protects test rows before model fitting.
    \item The training-mean baseline establishes minimum useful performance.
    \item \c{fit()} learns coefficients from training pairs only.
    \item \c{predict()} produces one numeric estimate per test row.
    \item MAE, $R^2$, and residual plots evaluate different aspects of error.
    \item Coefficients are printed with feature names so their units can be interpreted.
\end{enumerate}
Synthetic data prove that the code can recover a known pattern; they do not
validate performance on real buildings.
\end{SyntaxBox}

\hh{Understand coefficients without claiming causality}

With other features held fixed, a coefficient is the model's estimated change
in prediction for one unit increase in that feature. This interpretation assumes
the linear form is appropriate and features are measured consistently. Correlated
features can make individual coefficients unstable. An observational predictive
coefficient does not prove that changing the feature will cause the target to
change by that amount.

\hh{Recognize overfitting and leakage}

\textbf{Overfitting} occurs when a model learns training-specific detail that
does not generalize. Training performance may improve while test performance
worsens. \textbf{Leakage} occurs when training uses information unavailable at
prediction time or information derived from protected test outcomes.

Common leakage examples include:

\begin{itemize}
    \item filling missing values using means calculated from the full dataset;
    \item choosing features after comparing many options on the final test set;
    \item including a post-outcome inspection result as a predictor; and
    \item placing records from the same building in both training and test sets
          when deployment concerns entirely new buildings.
\end{itemize}

Use a scikit-learn \c{Pipeline} when preprocessing must learn parameters. Fit
the pipeline on training data so each transformation follows the same boundary
as the estimator.

\hh{Preview classification and clustering}

The estimator interface remains consistent across different algorithms.

\begin{lstlisting}[
    language=python,
    style=block,
    caption={Recognize the common estimator interface},
    label={c18_other_tasks},
    captionpos=b
]
from sklearn.cluster import KMeans
from sklearn.linear_model import LogisticRegression
from sklearn.metrics import confusion_matrix

# Classification: learn labeled categories.
classifier = LogisticRegression(max_iter=1000)
classifier.fit(X_train, complaint_train)
complaint_predictions = classifier.predict(X_test)
print(confusion_matrix(complaint_test, complaint_predictions))

# Clustering: discover feature-space groups without target labels.
clusterer = KMeans(n_clusters=3, random_state=42, n_init=10)
cluster_labels = clusterer.fit_predict(X)
\end{lstlisting}

Classification evaluation must examine which classes are confused, not only
overall accuracy. Clustering labels are arbitrary identifiers; a domain expert
must inspect whether groups are stable, useful, and ethically appropriate. The
official scikit-learn
\href{https://scikit-learn.org/stable/modules/generated/sklearn.cluster.KMeans.html}{KMeans reference}
documents the estimator's parameters and learned attributes.

\hh{Document intended use and limitations}

Every applied model should record:

\begin{itemize}
    \item the decision or prediction question and intended users;
    \item data source, dates, sample definition, exclusions, and missingness;
    \item features available at prediction time and the target definition;
    \item split strategy, baseline, metrics, and uncertainty;
    \item performance for important subgroups and failure cases;
    \item known distribution shifts and prohibited uses; and
    \item when retraining, monitoring, or human review is required.
\end{itemize}

Prediction is not explanation, and a high test score does not transfer
automatically to a different city, building type, season, sensor system, or
population.

\hh{Common mistakes and useful diagnoses}

\begin{itemize}
    \item \textbf{$X$ is one-dimensional.} Supervised estimators normally expect
          a table shaped \c{(samples, features)}.
    \item \textbf{Feature order changes.} Preserve DataFrame column names and
          order between fitting and prediction.
    \item \textbf{The test set is consulted repeatedly.} Use training and
          validation procedures for development; reserve the test for final evaluation.
    \item \textbf{Only $R^2$ is reported.} Include an error magnitude in target
          units, a baseline, and residual diagnostics.
    \item \textbf{Random splitting ignores structure.} Use time, group, or
          spatial splits when those relationships define deployment.
    \item \textbf{Association becomes a causal claim.} State predictive scope
          and avoid intervention language without a causal design.
\end{itemize}

\hh{Practice ladder}

\hhh{Level 0 -- Identify}

In the building example, identify the samples, four features, target, estimator,
baseline, and two evaluation metrics.

\hhh{Level 1 -- Modify}

Change \c{test_size} to 0.20 while keeping \c{random_state=42}. Report the new
train and test row counts and compare the metrics.

\hhh{Level 2 -- Explain}

Explain why calculating a missing-value mean from the full dataset before the
split is leakage, even though no target column is used.

\hhh{Level 3 -- Build}

Write \c{evaluate_regression(y_true, predictions, baseline_predictions)} that
returns a dictionary containing baseline MAE, model MAE, and $R^2$.

\hhh{Level 4 -- Challenge}

Create a grouped split in which all rows from one building identifier remain in
only training or test data. Compare the result with a random row split and
explain which deployment question each split represents.

\hh{Practice solutions}

\textbf{Level 0:} Rows are synthetic buildings; the features are floor area,
glazing ratio, occupants, and outdoor temperature; the target is energy; the
estimator is linear regression; the baseline is the training-target mean; the
metrics are MAE and $R^2$.

\textbf{Level 2:} The full-dataset mean contains information from the protected
test feature distribution. A preprocessing decision has therefore used data
that should have remained unseen.

\textbf{Level 3}

\begin{lstlisting}[
    language=python,
    style=block,
    caption={One reusable regression evaluation function},
    label={c18_solution_evaluate},
    captionpos=b
]
from sklearn.metrics import mean_absolute_error, r2_score

def evaluate_regression(y_true, predictions, baseline_predictions):
    return {
        "baseline_mae": mean_absolute_error(
            y_true,
            baseline_predictions,
        ),
        "model_mae": mean_absolute_error(y_true, predictions),
        "r2": r2_score(y_true, predictions),
    }
\end{lstlisting}

The function receives predictions rather than fitting a model, so evaluation
remains independent of the training algorithm.

\hh{Chapter checkpoint}

\begin{enumerate}
    \item How do regression, classification, and clustering differ?
    \item What shapes do $X$ and $y$ conventionally have?
    \item Why is a test set protected before fitting?
    \item Why should a baseline use training information only?
    \item What does MAE measure, and in what units?
    \item What can a residual plot reveal that one metric cannot?
    \item What is leakage, and how does it differ from overfitting?
    \item Why does predictive association not establish causality?
\end{enumerate}

\hh{Checkpoint answers}

\begin{enumerate}
    \item They predict numbers, predict categories, and discover unlabeled groups.
    \item $X$ is samples by features; $y$ contains one target per sample.
    \item To provide independent evidence about unseen-data performance.
    \item Otherwise test answers influence the comparison.
    \item Average absolute prediction error in the target's units.
    \item Systematic curves, changing spread, subgroups, and outliers.
    \item Leakage exposes unavailable or protected information; overfitting
          learns training detail that fails to generalize.
    \item Observed relationships may reflect confounding, selection, or measurement.
\end{enumerate}

\begin{tcolorbox}[title=Chapter summary,colframe=orange,colback=white,arc=2mm]
Machine learning estimates a rule from data, and the estimate is only as
trustworthy as the workflow around it: a legitimate target, a protected test
set, a baseline that uses training information only, metrics paired with
residual diagnostics, and stated limits on population and use.
\booklink{ch:19}{Chapter 19} opens the solvers that fit such models, and
\booklink{ch:20}{Chapter 20} extends this workflow to classification, where
several candidate models compete and the comparison itself must be protected
from the test set.
\end{tcolorbox}
